% !TEX program = lualatex
\RequirePackage{fix-cm}
\documentclass[a4paper,11pt]{ltjsarticle}
\usepackage{amsmath,amssymb,amsthm}
\usepackage{bm}

\theoremstyle{plain}
\newtheorem{lemma}{補題}
\newtheorem{theorem}{定理}

\title{構成的論理における $\neg\neg$-差分とアルゴリズム的ランダム性：\\実現可能性を介したChaitin不完全性定理の一般化}
\author{千葉将臣}
\date{2026-07-11}

\begin{document}
\maketitle

\begin{abstract}
直観主義論理において、意味論的には古典的に真（$\neg\neg A$）でありながら構成的には証明不可能な命題（$A$）の集合、すなわち「$\neg\neg$-差分」の構造は、これまで主に位相的・集合論的な観点から研究されてきた。本稿では、この論理的ギャップを、単一命題 $A$ と $\neg\neg A$ の差分そのものとしてではなく、停止判定に関する命題族 $\{E_N\}_N$ のうち構成的に到達不能となる添字の閾値現象として操作化する。我々は、Lawvere--Tierneyの $\neg\neg$-位相を用いた空間的アプローチ（疎集合と測度零集合の乖離）の限界を指摘し、Kleeneの実現可能性（realizability）と接頭辞Kolmogorov複雑度を直接結びつける代替的な枠組みを提案する。主定理として、直観主義算術において「長さ $N$ 以下の全プログラムの停止性を一様に決定する排中律」の実現にはChaitinの $\Omega$ と同等の情報量が要求されることを示し、公理系の複雑度 $K(F)$ を超える領域で $F \vdash \neg\neg E_N$ かつ $F \nvdash E_N$ が生じることを証明する。
\end{abstract}

\section{序論}

\subsection{背景と動機}

直観主義論理と古典論理の最大の違いは、排中律 $A \lor \neg A$ や二重否定の除去 $\neg\neg A \to A$ の扱いに存する。Glivenkoの定理やGödel--Gentzenの負解釈が示すように、古典論理は直観主義論理の中に「二重否定で翻訳された安定な命題のクラス（$\neg\neg$-層）」として埋め込むことができる。トポス理論においては、LawvereとTierney (1972) がこれを $\neg\neg$-位相（$\neg\neg$-topology）という概念で幾何学的に定式化し、任意のトポス内部に古典論理を満たすブール・トポスを構築する道を開いた。

本研究の核心的な動機は、「古典的な真理空間と直観主義的な構成空間の間に生じる『差分（$\neg\neg A \setminus A$）』とは、一体どのような数学的実体なのか」という問いである。直観主義的に $A$ が成立するためには、その命題を構成的に裏付ける証拠（実現子、realizer）が必要となる。したがって、「$\neg\neg A$ だが $A$ ではない」という状態は、対象が意味論的には確定しているにもかかわらず、そこへ到達するためのアルゴリズム的な構成ステップが欠如している状態を指す。

\subsection{位相的アプローチの限界と情報論的転回}

この「差分空間」を幾何学的・位相的に捉えようとする素朴な試みは、深刻な困難に直面する。Cantor空間上の位相的構成（Baire空間の性質等）を用いた場合、$\neg\neg$-差分はnowhere dense（疎集合、meager ideal）に対応する。一方で、アルゴリズム的情報理論（AIT）において「到達不能な情報」の代名詞であるMartin-Löfランダムな実数の集合は、Lebesgue測度零（null ideal）の補集合として定義される。Cichońの図式が示すように、meager idealとnull idealは独立した公理力を持つため、両者は素朴には一致しない。

この乖離は、現象を「外側からの空間の大きさ」で測ろうとしたことに起因する。本稿では、Hyland (1982) のEffective Topos（$\mathcal{Eff}$）の精神に則り、空間の形ではなく「実現子の構成に要求される計算複雑度（証明コスト）」という内側からの尺度を採用する。Schnorrの定理により、Martin-Löfランダム性は「Kolmogorov複雑度による非圧縮性」と完全に同値であることが知られている。我々はこの事実を利用し、位相的差分をそのまま同一視するのではなく、命題族 $\{E_N\}_N$ に対する到達不能性として再定式化することで、直観的推論の限界とAITにおける情報論的限界（Chaitinの不完全性定理）を比較可能な形に置き直す。

なお、本稿の $D_F(E_N)$ は、もとの問題意識である「単一の命題 $A$ と $\neg\neg A$ の間に何があるか」をそのまま集合として取り出したものではない。むしろ、停止判定排中律の有限近似 $E_N$ をパラメータ化し、どの有限段階から先で $F$ が構成的実現子を供給できなくなるかを測る指標である。この操作化により、$\neg\neg$-差分を「命題族に沿った閾値現象」として扱う。

\subsection{関連研究と本稿の位置づけ}

Chaitinの不完全性定理は、十分に強い健全な形式体系が、その記述複雑度を大きく超えるKolmogorov複雑度をもつ具体的対象の非圧縮性を証明できないことを示す。また、Solovayによる停止確率 $\Omega$ の研究は、$\Omega$ の初期ビットが有限長プログラムの停止情報を符号化することを明確にした。したがって、本稿で用いる情報論的下界そのものは、AITにおける既存の標準的事実に基づいている。

本稿の新規性は、これらの既存結果を直観主義的証明論の言葉へ移し替える点にある。具体的には、Kleene実現可能性とNelsonの抽出定理を明示的に経由することで、$F \vdash E_N$ という構成的証明可能性から一様停止判定の実現子が抽出され、その実現子の複雑度が $K(F)+O(\log N)$ に抑えられることを示す。この上界と、$\Omega$ 由来の下界 $N-O(1)$ を衝突させることで、Chaitin型の情報論的不完全性を「$\neg\neg E_N$ は証明できるが $E_N$ は証明できない」というrealizability上の到達不能性として表現する。

\section{予備知識}

本稿で用いるアルゴリズム的情報理論、および直観主義的証明論の基本概念を定義する。

\subsection{接頭辞Kolmogorov複雑度とChaitinの $\Omega$}

チューリング完全な\textbf{万能接頭辞機械（universal prefix-free machine）} $U$ を一つ固定する。入力プログラム $p \in \{0, 1\}^*$ は接頭辞符号（ある有効なプログラムが別の有効なプログラムの接頭辞にならない）をなすとする。

対象となる文字列または自然数 $x$ に対する\textbf{接頭辞Kolmogorov複雑度} $K(x)$ を、以下のように定義する。
\[
K(x) := \min \{ \lvert p\rvert : U(p) = x \}.
\]
ここで $\lvert p\rvert$ はプログラム $p$ のビット長を表す。また、万能機械 $U$ が停止する確率（Chaitinの定数）を以下で定義する。
\[
\Omega := \sum_{U(p)\downarrow} 2^{-\lvert p\rvert}.
\]
$\Omega$ は計算不可能な実数であり、その二進展開の最初の $N$ ビットを $\Omega_N$ としたとき、$K(\Omega_N) \ge N - O(1)$ を満たす（すなわち、漸近的に非圧縮である）ことが知られている。

\subsection{直観主義算術とKleeneの実現可能性}

基礎となる形式体系として、\textbf{直観主義算術（Heyting Arithmetic, HA）}、あるいはそれを拡張した健全かつ帰納的に公理化可能な体系 $F$ を対象とする。体系 $F$ の公理と推論規則を列挙するプログラムの最小ビット長を $K(F)$ と記述する。

Kleeneのnumerical realizability関係 $e \Vdash A$（自然数 $e$ が論理式 $A$ を実現する）を標準的に定義する。特に以下の性質に留意する。

\begin{itemize}
\item \textbf{停止問題の述語:} $T(n, n, s)$ を「プログラム $n$ が入力 $n$ に対して $s$ ステップで停止する」という原始帰納的関係とする。停止問題は $H(n) \equiv \exists s\, T(n, n, s)$ と表される。
\item $e \Vdash \exists s\, T(n, n, s)$ とは、$e$ を計算すると具体的な停止ステップ数 $s$ を出力し、かつ $T(n, n, s)$ が真であることを意味する。
\item $e \Vdash (A \lor B)$ とは、$e$ を計算するとペア $\langle i, p \rangle$ が得られ、$i=0$ なら $p \Vdash A$、$i=1$ なら $p \Vdash B$ となることである。
\end{itemize}

さらに、構成的体系の根幹をなす事実として\textbf{Nelsonの抽出定理}（1947）を用いる。すなわち、$F \vdash A$ であるならば、その証明のGödel数から、アルゴリズム的に実際の実現子 $e$（$e \Vdash A$）を抽出する部分再帰的関数が存在する。


\section{枠組みの定式化}

\subsection{一様停止判定命題と情報量}

本稿では、対象とする命題として、単一のプログラムの停止問題ではなく、「一定の長さ以下のすべてのプログラムの停止性を一様に決定する命題」を採用する。

万能接頭辞機械 $U$ への入力プログラム $p$ に関する停止問題を $H(p) \equiv \exists s\, T_U(p, p, s)$ と表す。このとき、長さ $N$ 以下のすべてのプログラムに関する排中律 $E_N$ を以下のように定義する。
\[
E_N \equiv \forall p \in \{0, 1\}^{\le N}\, (H(p) \lor \neg H(p)).
\]
ここで、最初の $N$ 個の自然数に対する停止性ではなく、「長さ $N$ 以下の接頭辞プログラム」を対象とする点が決定的に重要である。最初の $N$ 個のプログラムのうち停止するものの個数（高々 $N$）を指定するための情報量は $\log_2 N$ ビットで済むが、長さ $N$ 以下の全接頭辞プログラム（最大 $2^{N+1}-1$ 個）の停止性を網羅的に決定するためには、本質的にChaitinの定数 $\Omega$ の最初の $N$ ビットに相当する情報量が要求される。

\subsection{$\neg\neg$-差分集合の定義}

健全で帰納的に公理化可能な直観主義形式体系 $F$（直観主義算術HAを含む）を固定する。体系 $F$ に相対的な古典的真理と構成的真理のギャップを測るため、以下の$\neg\neg$-差分集合を定義する。
\[
D_F(E_N) := \{ N \in \mathbb{N} : F \vdash \neg\neg E_N \text{ かつ } F \nvdash E_N \}.
\]
古典論理のモデル（Lawvere--Tierneyの $\neg\neg$-層）においては、排中律の二重否定は自明に真であるため、この集合は「意味論的には確定しているが、体系 $F$ 内部での構成的証明が存在しない添字 $N$ の集合」を正確に表現している。本研究の目的は、この論理的な差分集合が空でないことを示すだけでなく、その発生境界がアルゴリズム的情報理論の限界と正確に一致することを証明することである。

\section{主定理とその証明}

本章では、体系 $F$ の複雑度 $K(F)$ と、$E_N$ の構成的証明から抽出される実現子の複雑度を評価することで、主定理を導く。

\subsection{証明からの実現子抽出}

最初の補題として、体系 $F$ において $E_N$ が証明可能であると仮定した場合の、実現子の情報論的な上限を定める。

\begin{lemma}[証明複雑度による実現子の上界]
健全で帰納的に公理化可能な直観主義形式体系 $F$ において、$F \vdash E_N$ が成立するならば、$E_N$ を実現するあるインデックス $e \in \mathbb{N}$（すなわち $e \Vdash E_N$）が存在し、その接頭辞Kolmogorov複雑度 $K(e)$ は以下の上界を満たす。
\[
K(e) \le K(F) + O(\log N).
\]
ここで隠れた定数は $N$ および $F$ に依存しない。
\end{lemma}

\begin{proof}
$F \vdash E_N$ を仮定する。このとき、万能機械 $U$ 上で動作する以下のアルゴリズム $\mathcal{M}$ を構成する。
\begin{enumerate}
\item \textbf{証明の探索:} 体系 $F$ の定理を列挙するプログラム（サイズ $K(F)$）と、自然数 $N$ の自己区切り表現（サイズ $O(\log N)$）を用いて、$F$ の定理を順次生成し、結論が $E_N$ のGödel数と一致する証明 $p_{E_N}$ を発見する。仮定より、この探索は有限時間で停止する。
\item \textbf{実現子の抽出:} Nelsonの抽出定理（1947）により、構成的体系の証明 $p_{E_N}$ から、その結論の実現子を抽出する部分再帰的関数 $\varepsilon$ が存在する。$\mathcal{M}$ は $\varepsilon(p_{E_N})$ を計算し、その結果 $e$ を出力する。
\end{enumerate}
アルゴリズム $\mathcal{M}$ 自身の動作コードと関数 $\varepsilon$ の記述長は定数 $O(1)$ であるため、万能機械 $U$ が $e$ を出力するために必要な全プログラムサイズは
\[
K(e) \le K(F) + K(N) + O(1) \le K(F) + O(\log N)
\]
となる。
\end{proof}

\subsection{実現子の情報論的下界}

次に、$E_N$ を実現するいかなるアルゴリズムも、本質的にChaitinの $\Omega$ と同等の情報量を内包せざるを得ないことを示す。

\begin{lemma}[ランダム性による実現子の下界]
命題 $E_N$ を実現する任意のインデックス $e$（$e \Vdash E_N$）について、その接頭辞Kolmogorov複雑度は以下の下界を満たす。
\[
K(e) \ge N - O(1).
\]
ここで隠れた定数は $N$ に依存しない。
\end{lemma}

\begin{proof}
$e \Vdash E_N$ と仮定する。実現可能性の定義より、$e$ は長さ $N$ 以下の任意のプログラム $p$ を入力として受け取り、その停止性の真理値フラグを正しく出力する計算可能関数として機能する。この $e$ をサブルーチンとして用いることで、長さ $N$ 以下の全プログラムの停止性を判定し、以下の有限和を正確に計算できる。
\[
\Omega^{(N)} = \sum_{\lvert p\rvert \le N,\, U(p)\downarrow} 2^{-\lvert p\rvert}.
\]
Chaitin（1975）の定理により、長さ $N$ 以下の全停止プログラムの集合から、真の停止確率 $\Omega$ の二進展開の最初の $N$ ビット $\Omega_N$ を抽出する計算可能関数が存在する。したがって、$e$ に $O(1)$ の定数長プログラムを結合するだけで $\Omega_N$ を出力可能である。
\[
K(\Omega_N) \le K(e) + O(1).
\]
一方で、$\Omega$ はMartin-Löfランダムな実数であり、AITの基本定理により $K(\Omega_N) \ge N - c$ が成り立つ。これらを結合し、$K(e) \ge N - O(1)$ を得る。
\end{proof}

\subsection{差分の発生境界}

以上の準備の下で、論理の差分がAITの非圧縮性によって生じることを決定的に証明する。

\begin{theorem}[$\neg\neg$-差分の情報論的境界]
任意の健全な直観主義形式体系 $F$ に対して、ある定数 $c, d > 0$ が存在し、次が成り立つ。
\[
N > K(F) + c \log N + d \implies N \in D_F(E_N).
\]
すなわち、体系の複雑度 $K(F)$ を漸近的に超えるすべての $N$ は、古典的真理と構成的真理の差分空間に属する。
\end{theorem}

\begin{proof}
条件を満たす $N$ について、$N \in D_F(E_N)$、すなわち (1) $F \vdash \neg\neg E_N$ かつ (2) $F \nvdash E_N$ の両方を示す。

\textbf{Part 1: 古典的自明性}
直観主義論理において $\neg\neg(A \lor \neg A)$ は定理である。$E_N$ は有限集合 $\{0, 1\}^{\le N}$ 上の量化であり、有限個の論理積に展開可能である。直観主義論理において二重否定は有限の論理積に対して分配的であるため、任意の $N$ について $F \vdash \neg\neg E_N$ が常に成立する。

\textbf{Part 2: 構成的不到達性}
背理法を用いる。ある $N$ について $F \vdash E_N$ が成立したと仮定する。補題1より、$E_N$ を実現する $e$ が存在し、$K(e) \le K(F) + c_1 \log N + c_2$ を満たす。一方、補題2より同じ $e$ は $K(e) \ge N - c_3$ を満たさねばならない。結合すると、
\[
N - c_3 \le K(F) + c_1 \log N + c_2
\]
であり、したがって
\[
N \le K(F) + c_1 \log N + (c_2 + c_3)
\]
となる。しかし、定数 $c = c_1$、$d = c_2 + c_3$ と置いたとき、$N > K(F) + c \log N + d$ となる $N$ を選択すれば論理的矛盾が生じる。体系 $F$ の健全性により仮定は誤りであり、この領域では $F \nvdash E_N$ でなければならない。

以上より、境界を超えたすべての $N$ において $F \vdash \neg\neg E_N$ かつ $F \nvdash E_N$ であり、$N \in D_F(E_N)$ が証明された。
\end{proof}



\section{トポス理論的解釈：$\mathcal{Eff}$ と $\neg\neg$-位相の幾何学}

前章で示された主定理は、証明論的・計算量的な限界を与える定理である。本章では、それをHylandのEffective Topos $\mathcal{Eff}$ における幾何学的・位相的な構造と対応づける解釈を述べる。ただし以下の議論は、主定理から直ちに従う証明済みの主張と、それに基づくトポス論的な読み替えとを区別して提示する。

\subsection{$\mathcal{Eff}$ における実数対象と $\Omega$ の位置づけ}

$\mathcal{Eff}$ においては、実数の定義（Cauchy実数、Dedekind実数、MacNeille実数など）が、古典的な集合論のモデル（$Set$）とは異なり一致しないことが知られている。Chaitinの定数 $\Omega$ は、その定義から左切断（有理数の下界）が帰納的可枚挙であるため、$\mathcal{Eff}$ 内部において下半実数（lower real）としては正当な対象として存在する。

しかし、$\Omega$ が $\mathcal{Eff}$ において完全な実数（Dedekind実数）として扱われるためには、十分細かい精度で有理近似との大小関係を決定する実現子が必要になる。この要求は、有限精度では命題 $E_N$ を実現するアルゴリズムの要求と同型の情報を含む。

\begin{theorem}[有限精度実数化の障害]
任意の健全な直観主義形式体系 $F$ に対して、ある定数 $c,d>0$ が存在し、$N>K(F)+c\log N+d$ ならば、$F$ は $\Omega$ の最初の $N$ ビットを決定する実現子を与える証明を持たない。
\end{theorem}

\begin{proof}
もしそのような実現子を $F$ の証明から抽出できるならば、$\Omega_N$ から長さ $N$ 以下の接頭辞プログラムの停止性を復元する標準的な計算を経由して、$E_N$ の実現子が得られる。これは主定理の $F \nvdash E_N$ と矛盾する。
\end{proof}

したがって、「$N \approx K(F)$ の段階で $\Omega$ の実数としての実体化が破綻する」という表現は、厳密には上の有限精度版の系をトポス論的に読んだものである。すなわち、$\Omega$ は $\mathcal{Eff}$ 内部において下半実数としては自然に現れるが、体系 $F$ が供給できる実現子だけでは、十分大きな精度でDedekind実数的な決定性を与えられない対象として振る舞う。

\subsection{$\neg\neg$-層化によるオラクル注入と空間の差分}

LawvereとTierneyの $\neg\neg$-位相は、$\mathcal{Eff}$ 内部に古典論理を満たすブール・部分トポス（$\neg\neg$-層）を切り出す。この $\neg\neg$-層は、外部の古典的な集合界 $Set$ と同値になることが知られている。

本稿の枠組みにおいて、$\neg\neg$ 演算子による二重否定シフトは、単なる論理規則ではなく、実現子の構成プロセスを保留し、外部の古典的決定性を参照する操作として解釈できる。主定理が示す $N \in D_F(E_N)$ という差分領域の発生は、この意味で「$\mathcal{Eff}$ の構成的世界では実現子が得られないが、$\neg\neg$-層では古典的に安定化される」有限段階を特定している。

これまで、Cantor空間におけるこの種の「差分」を位相的に捉えると、meager ideal（疎集合）とnull ideal（測度零）の不一致に直面する。本稿はこの不一致を同一空間上で解消したわけではない。むしろ、位相的定義を採用し続けるのではなく、命題族 $\{E_N\}_N$ の添字集合 $D_F(E_N)$ へと問題を移し、トポス内部での\textbf{証明コスト（接頭辞Kolmogorov複雑度）}によって代替的に定式化した。この定式化は、当初の動機である「構成的証拠の欠如」をより直接に測るものであり、オラクルなしでは圧縮不可能な情報が、どの有限段階から構成的到達不能性として現れるかを記述する。

\section{結論と今後の展望}

\subsection{結論}

本研究は、「直観主義論理と古典論理の差分」という問題意識を、停止判定排中律の命題族 $\{E_N\}_N$ に沿った閾値現象として操作化した。我々は、Lawvere--Tierneyの $\neg\neg$-位相を用いた空間的・測度的なアプローチが抱える齟齬を回避し、Kleeneの実現可能性と接頭辞Kolmogorov複雑度を直接結びつける代替的な枠組みを構築した。

この枠組みにおいて、直観主義算術における一様停止判定命題 $E_N$ の構成的証明には、Chaitinの $\Omega$ と等価な情報量が不可避に要求される（補題2）。一方で、形式体系 $F$ が生成できる実現子の情報量は $K(F)$ にバウンドされる（補題1）。この証明論的限界と情報論的限界の衝突から導かれる主定理により、公理系の複雑度を漸近的に超える領域（$N > K(F) + c \log N + d$）において、古典的には自明な命題が直観主義的には完全に到達不能になることが示された。

この意味で、$D_F(E_N)$ に現れる $\neg\neg$-差分は、体系から見て予測不可能なランダムビットが構成的到達不能性として現れる場所を示している。本結果は、Chaitinの不完全性定理をrealizability toposの観点から読み替え、論理学・計算機科学・圏論の結びつきを明確化するものである。

\subsection{今後の展望}

本研究の成果は、証明論とアルゴリズム的情報理論の交差点において、さらに多くの研究の方向性を示唆している。

\begin{enumerate}
\item \textbf{他のランダム性概念への拡張:} 本稿ではMartin-Löfランダム性（Chaitinの $\Omega$）と接頭辞複雑度に焦点を当てたが、Schnorrランダム性や計算可能ランダム性など、他のランダム性の階層が $\mathcal{Eff}$ のどのような論理的断片、あるいは別の位相・closure operatorと対応するかを解明することは興味深い課題である。
\item \textbf{直観主義的逆数学との接続:} $E_N$ に代表される命題群を、Markovの原理（Markov's Principle, MP）やLPO（Limited Principle of Omniscience）といった構成的数学の諸原理と複雑度論的に関連付けることで、公理系の強さを「ランダムビットの抽出能力」として測り直す新たな逆数学の構築が期待される。
\item \textbf{より高次のトポスにおける不完全性:} 本枠組みを、算術に留まらず高階論理や型理論に対応する実現可能性トポス（例：modified realizability topos）へと拡張し、計算の停止性以外の「高次の情報論的非圧縮性」がどのような論理的差分を生み出すかを幾何学的に研究することが次のステップとなろう。
\end{enumerate}

\end{document}
